\section{Gaussian decomposition of the entropy deficit}
\label{mat:section}

The transport estimate of Proposition~\ref{tr:surplus} controls the
Gaussian energy of the reparametrized Jacobian.  Its degree-one part,
however, has no Poincar\'e surplus.  We use that part to choose a matrix
supporting plane for the log determinant.  The remaining Gaussian
chaoses then supply the coercivity needed to prove the entropy bound.
We retain a strict remainder that will also identify the equality cases.
Throughout this section, $\E$ denotes integration against the standard
Gaussian measure $\gamma_m$.

Define a scalar weight on $[0,\infty)$ by
\begin{equation}\label{mat:rigidity-weight}
 \omega(\lambda)=
 \frac{3(\lambda-1)^2}{37+3\lambda}
 \left(\frac3{20}+\frac{49}{100}\frac{37}{37+3\lambda}\right)
 \min\left\{\frac{13}{100},\frac{32}{25\sqrt\lambda}\right\},
\end{equation}
where the minimum is $13/100$ at $\lambda=0$.
This weight is continuous and nonnegative, vanishes only at $1$, and
satisfies $\omega(\lambda)/\sqrt\lambda\to24/125$ as
$\lambda\to\infty$.

\begin{theorem}[The smooth entropy bound with a remainder]
\label{mat:entropy-smooth}\label{prop:slack}
Let $m\ge1$ be an integer and let $\mu=e^{-V}\,\dd y$ be a
probability measure on $\R^m$, where
$V\in C^\infty(\R^m)$ and
$c_-\Id\le D^2V\le c_+\Id$ for constants $0<c_-\le c_+<\infty$.
Choose the affine normalization of Proposition~\ref{tr:normalization},
and let $A=q^{-1}(J)$ be the resulting Gaussian transport matrix field,
so $\E A=0$.  Set
\[
 X_r=\E(x_rA),\qquad Y=A-\sum_{r=1}^m x_rX_r,\qquad
 B=\sum_{r=1}^mX_r^2.
\]
If $\lambda_1,\ldots,\lambda_m$ are the eigenvalues of $B$, then
\begin{equation}\label{mat:strict-remainder}
 2h(\mu)-2m-\log\det\Cov(\mu)
 \ge \frac1{500}\E\tr(Y^2)
       +\frac1{1000}\sum_{i=1}^m\omega(\lambda_i).
\end{equation}
In particular,
\[
 h(\mu)\ge m+\frac12\log\det\Cov(\mu).
\]
\end{theorem}

We use the affine normalization of
Proposition~\ref{tr:normalization}: the Brenier Jacobian $J$ and the
symmetric matrix field $A=q^{-1}(J)$ satisfy $\E A=0$.  This normalization
does not change $h(\mu)-\frac12\log\det\Cov(\mu)$.  For the proof,
we replace $\mu$ by this affine image and retain the notation $\mu$.
Scalar functions of $A$
are interpreted by spectral calculus; when their arguments are omitted,
$q,k,t,b,l$ and the functions introduced below are evaluated at $A$.  For a matrix field $F$, put
\[
 \tau(F)=\E\tr F,\qquad
 \|F\|^2=\E\tr(F^tF),\qquad
 \|F\|_S^2=\E\tr(F^tSF)
\]
when $S$ is a constant positive definite matrix.  For a family indexed
by $r$, these squared norms include summation over $r$.  Lemma~\ref{tr:energy}
provides the Sobolev regularity used below.

Write $C=\Cov(\mu)$ and let
\[
 \mathscr D=2m+\log\det C-2h(\mu)
\]
be the negative of twice the entropy gap.  The transport change of variables and
$\log q(a)=-a^2/2-\frac12\log(2\pi)+t(a)$ give
\begin{equation}\label{mat:deficit-exact}
 \mathscr D=\tau(A^2-2t)+\log\det C+m.
\end{equation}

\subsection{Covariance and the first Gaussian chaos}

Let $N=-\Delta+x\cdot\nabla$ be the nonnegative Ornstein--Uhlenbeck
operator and let $R=(\Id+N)^{-1}$.  The degree-$d$ Gaussian chaos is
the span of coordinatewise products of one-dimensional Hermite polynomials whose
degrees sum to $d$.  These spaces give an orthogonal decomposition of
$L^2(\gamma_m)$; on the degree-$d$ space, $N$ and $R$ act by $d$ and
$(1+d)^{-1}$, respectively.
We write $\partial=(\partial_1,\ldots,\partial_m)$ and
$\partial^*v=\sum_r(x_rv_r-\partial_rv_r)$, so $N=\partial^*\partial$.
The usual Gaussian Sobolev norm is
$\|F\|_{H^1}^2=\|F\|^2+\|\partial F\|^2$; its dual norm is
\[
 \|F\|_{-1}^2=\sum_{d\ge0}\frac{\|F_d\|^2}{1+d},
\]
initially for square-integrable fields, and then by completion.
Here $F_d$ denotes the degree-$d$ chaos component.  These spectral
conventions are standard; see \cite[Section~2]{NourdinPeccati2009}.

We use the Gaussian Helffer--Sj\"ostrand covariance representation;
see \cite[Proposition~4.1(ii)]{DuerinckxOtto2020}.  We include a direct
chaos proof below.  The term $-2\tau(t)$ in \eqref{mat:deficit-exact}
determines its completion.  Since $qk=t$, subtracting $(\Id+N)k$ from
$q$ in the covariance quadratic form produces exactly the cross term
$-2t$.

\begin{lemma}[Covariance completion]\label{mat:covariance}
Define the symmetric $H^{-1}$ matrix field
\[
 U=q(A)-(\Id+N)k(A)
\]
and let $W$ be the Gram matrix of its rows in $H^{-1}$:
\[
 W_{ij}=\sum_{\ell=1}^m
       \langle U_{i\ell},U_{j\ell}\rangle_{-1}.
\]
Then
\begin{equation}\label{mat:covariance-completion}
 C=\E\bigl[J(RJ)\bigr]
   =W+\E\left(2t-k^2-\sum_r K_r^2\right),
 \qquad K_r=\partial_r k(A).
\end{equation}
Moreover, if $U=u-\partial^*v$ with $u$ and the family $v=(v_r)_r$
square-integrable, then for every constant $S>0$,
\begin{equation}\label{mat:dual-bound}
 \tr(SW)\le\|u\|_S^2+\|v\|_S^2.
\end{equation}
\end{lemma}

\begin{proof}
Put $F=\nabla\phi$.  For two degree-$d$ components of $F$, Gaussian
integration by parts gives
\[
 \sum_\ell\E[(\partial_\ell F_{i,d})(\partial_\ell F_{j,d})]
 =d\,\E[F_{i,d}F_{j,d}].
\]
Differentiation lowers the degree by one, and $R$ therefore cancels the
factor $d$.  Summing over $d\ge1$ proves
$C_{ij}=\sum_\ell\E[J_{i\ell}R J_{j\ell}]$.  Symmetry of $J$ gives
the first matrix identity in \eqref{mat:covariance-completion}.
Expanding the row Gram matrix of $q-(\Id+N)k$ gives cross terms
$-2\E(qk)=-2\E t$.  The remaining term is
\[
 \E\bigl[k(\Id+N)k\bigr]
   =\E\left(k^2+\sum_r K_r^2\right),
\]
by entrywise integration by parts.  This proves the second identity.
The calculation is valid by the $H^1$--$H^{-1}$ pairing, even if $Nk$
is not square-integrable.

For a test field $F$, the pairing with $u-\partial^*v$ is
$\langle u,F\rangle-\langle v,\partial F\rangle$.  Cauchy--Schwarz
bounds it by
$(\|u\|^2+\|v\|^2)^{1/2}\|F\|_{H^1}$.  Taking the dual norm proves
\eqref{mat:dual-bound} for $S=\Id$, and applying this argument to
$S^{1/2}U$ proves the general case.
\end{proof}

Separate the degree-one part of $A$ by writing
\begin{equation}\label{mat:chaos}
 A=\sum_{r=1}^m x_rX_r+Y,\qquad
 X_r=\E(x_rA),\qquad B=\sum_rX_r^2.
\end{equation}
The matrices $X_r$ are constant and symmetric; $Y$ contains only chaoses
of degree at least two.  We shall use
\begin{equation}\label{mat:gradient-fields}
 Z_r=\partial_rA=X_r+G_r,\qquad
 G_r=\partial_rY,\qquad
 g_r=\partial_rN^{-1}Y,
 \qquad V_0=\|G\|^2-\|Y\|^2.
\end{equation}
Thus $\partial^*g=Y$.  All three derivative fields are symmetric.

For comparison, the coordinatewise Gaussian-to-exponential transport
satisfies
\[
 A(x)=\operatorname{diag}(x_1,\ldots,x_m),\qquad B=\Id.
\]
We therefore choose the supporting plane for $\log\det C$ as a
regularized inverse of $B$ centered at this value.  Fix
\begin{equation}\label{mat:dual-constants}
 \alpha=\frac{3}{40},\qquad s=\frac1{1-\alpha},
\end{equation}
and set
\begin{equation}\label{mat:dual-matrices}
 S=(\Id/s+\alpha B)^{-1},\qquad T=S-\Id,
 \qquad H=-T(B-\Id).
\end{equation}
These constant matrices commute with one another, though generally not
with functions of $A$.  Directly from their definition,
\begin{equation}\label{mat:dual-relations}
 0<S\le s\Id,\qquad T=-\alpha S(B-\Id),\qquad
 H=\alpha S(B-\Id)^2\ge0,\qquad T^2=\alpha HS.
\end{equation}
In particular, $S=\Id$ when $B=\Id$.  Deviations of the first chaos from
this value produce the nonnegative matrix $H$.

Concavity of the log determinant yields
$\log\det C+m\le\tr(SC)-\log\det S$.  Combining this with
Lemma~\ref{mat:covariance} and Proposition~\ref{tr:surplus}, and using
$\|A\|^2-\|Z\|^2=-V_0$, gives
\begin{equation}\label{mat:deficit-reduction}
 \mathscr D\le -V_0-\mathcal C+\tr(T-\log S)+\tr(SW)
   +\tau\left[T\left(2t-k^2-\Id-\sum_rK_r^2\right)\right].
\end{equation}
Here $\mathcal C$ is the eigenvalue-separation surplus in
Proposition~\ref{tr:surplus}.  The remaining task is to estimate the last
two terms while retaining the negative contribution involving $H$.

\subsection{An exact expansion with noncommuting factors}

Recall the divided differences from Section~\ref{tr:section}, with
$v[a,a]=v'(a)$.  For a smooth scalar function $v$, define its secant
increment at the second endpoint by
\[
 \Delta v_{ij}=v[a_i,a_j]-v'(a_j).
\]
We use the matrix derivative formula in
\cite[Section~2.4]{Noferini2017}.
At a fixed point, choose an orthonormal eigenbasis of $A$.  Entrywise
multiplication is denoted by $\circ$.  Define the derivative residuals
\begin{equation}\label{mat:residuals}
 D_r=K_r-X_rb=G_rb+Z_r\circ\Delta k,\qquad
 L_r=\partial_rb-X_rl=G_rl+Z_r\circ\Delta b.
\end{equation}
These matrices need not be symmetric.  The distinction between $D_r$
and $D_r^t$ will be essential below.

The scalar identity $q-k=ab-l$ implies
\begin{equation}\label{mat:U-representation}
 U=Yb+\sum_rX_rL_r+(B-\Id)l-\partial^*D.
\end{equation}
Indeed,
$\partial^*(Xb)=(A-Y)b-Bl-\sum_rX_rL_r$.
To divide the higher-chaos terms between the two squares, fix
\begin{equation}\label{mat:constants}
 \beta=\frac7{25},\qquad c=\frac32,\qquad \gamma=\frac c2.
\end{equation}
Using $\partial^*g=Y$ in \eqref{mat:dual-bound}, we obtain
\begin{equation}\label{mat:W-bound}
 \tr(SW)\le
 \left\|Y(b-\beta)+\sum_rX_rL_r+(B-\Id)l\right\|_S^2
 +\|D-\beta g\|_S^2.
\end{equation}

We group the expansion according to its $H$ term and its derivative
residuals.  In the remaining $Y$ coefficient, we
split off $ck$: the identity $\partial^*g=Y$ will transfer
$c\tau(TYk)$ to derivatives, where completing the square produces
$D-\gamma g$.  Introduce the scalar functions
\begin{equation}\label{mat:scalar-functions}
 \begin{aligned}
 f&=k(1+b),&
 h_0&=f-ck-\frac{2(b-\beta)l}{\alpha},\\
 p&=f'-b^2-\frac{l^2}{\alpha},&
 e&=f'-2b^2-\frac{2l^2}{\alpha}.
 \end{aligned}
\end{equation}
The identity
\[
 af-f'=2t-k^2-b^2-1
\]
follows from $q-k=ab-l$ and $b=1-k(q-a)$.
For each pair of eigenvalues of $A$, put
\[
 \psi_{ij}=\Delta f_{ij}-2b_j\Delta k_{ij}
                         -\frac{2l_j}{\alpha}\Delta b_{ij},
 \qquad b_j=b(a_j),\quad l_j=l(a_j),
\]
and define
\begin{equation}\label{mat:M-w}
 M_r=G_re+Z_r\circ\psi+cg_rb,\qquad w_r=D_r-\gamma g_r.
\end{equation}
Finally, write
\begin{equation}\label{mat:P-definition}
 P_T=\tau(TbBb-TBb^2).
\end{equation}

\begin{lemma}[Matrix expansion]\label{mat:expansion}
The sum of the last two terms in \eqref{mat:deficit-reduction} is at
most
\begin{align}
 &-\tau(Hp)+\tau\left[T\left(Yh_0+\sum_rX_rM_r\right)\right]
       \label{mat:expanded-bound}\\
 &\quad+\left\|Y(b-\beta)+\sum_rX_rL_r\right\|_S^2
       +\|D-\beta g\|_S^2+\gamma^2\tau\left(T\sum_rg_r^2\right)\nonumber\\
 &\quad-P_T-2\tau\left([T,b]\sum_rX_rw_r\right)
                 -\tau\left(T\sum_rw_r^tw_r\right).\nonumber
\end{align}
Except for the use of \eqref{mat:W-bound}, this expansion is exact.
\end{lemma}

\begin{proof}
Let $F_0=Y(b-\beta)+\sum_rX_rL_r$.  The relations
\eqref{mat:dual-relations} give
\[
 \|F_0+(B-\Id)l\|_S^2
 =\|F_0\|_S^2+\frac1\alpha\tau(Hl^2)
                    -\frac2\alpha\tau(TF_0l).
\]
For the cross term, the order of the factors follows from
$\tr(F_0^tTl)=\tr(TF_0l)$, by transposition and cyclicity.
Integration by parts on the linear part of $A$ yields
\[
 \tau(TAf)=\tau\left[T\left(Yf+\sum_rX_r\partial_rf\right)\right],
 \qquad
 \partial_rf=X_rf'+G_rf'+Z_r\circ\Delta f.
\]
Together with $af-f'=2t-k^2-b^2-1$, this accounts for the terms involving
$f'$ and $f$.

To expand the energy term, use both versions of
$K_r=X_rb+D_r=bX_r+D_r^t$:
\begin{equation}\label{mat:K-square}
 \tau\left(T\sum_rK_r^2\right)
 =\tau(TbBb)+2\tau\left(Tb\sum_rX_rD_r\right)
             +\tau\left(T\sum_rD_r^tD_r\right).
\end{equation}
The two cross terms agree under transposition inside the trace.  Moving
$b$ past $T$, rather than past $X_r$, gives
\[
 \tau(TbX_rD_r)=\tau(TX_rD_rb)+\tau([T,b]X_rD_r).
\]
The terms containing $B-\Id$ now combine as
\[
 \tau\bigl(T(B-\Id)f'\bigr)+\tau(Tb^2-TbBb)
                 +\frac1\alpha\tau(Hl^2)
 =-\tau(Hp)-P_T.
\]
In the other terms, substituting \eqref{mat:residuals} produces
$G_re+Z_r\circ\psi$ and the $Y$ multiplier
$f-2(b-\beta)l/\alpha=h_0+ck$.

It remains to handle $c\tau(TYk)$.  Since $\partial^*g=Y$,
\[
 \tau(TYk)=\sum_r\tau(TK_rg_r).
\]
Here integration by parts first gives $\tau(Tg_rK_r)$; transposition
makes the two expressions equal.  Substituting $K_r=bX_r+D_r^t$ and
moving $b$ past $T$ introduces $cg_rb$ in $M_r$ and replaces $D_r$ by
$D_r-\gamma g_r$ in the commutator term.  The remaining square is
completed by
\[
 -\tau(TD_r^tD_r)+c\tau(TD_r^tg_r)
 =-\tau(Tw_r^tw_r)+\gamma^2\tau(Tg_r^2).
\]
This proves \eqref{mat:expanded-bound}.
\end{proof}

\subsection{The commutator and the weighted squares}

The commutator in Lemma~\ref{mat:expansion} has a compensating quadratic
term $P_T$.  Keeping that term is what permits a dimension-free estimate
without assuming that the Jacobian and the first-chaos matrix commute.

\begin{lemma}[Commutator estimate]\label{mat:commutator}
For the matrices $X_r,B,S,T$ in \eqref{mat:chaos} and
\eqref{mat:dual-matrices}, any square-integrable symmetric matrix field $b$, and arbitrary
square-integrable matrix fields $w_r$,
\begin{equation}\label{mat:commutator-bound}
 -P_T-2\tau\left([T,b]\sum_rX_rw_r\right)
       -\tau\left(T\sum_rw_r^tw_r\right)
 \le\left(1+\frac s8\right)\|w\|^2.
\end{equation}
\end{lemma}

\begin{proof}
Work pointwise and diagonalize $S$.  The relation
$B=\alpha^{-1}(S^{-1}-\Id/s)$ gives
\begin{equation}\label{mat:P-square}
 P_T=\frac1{2\alpha}
       \bigl\|S^{-1/2}[T,b]S^{-1/2}\bigr\|^2.
\end{equation}
Indeed, the contribution of $b_{ij}^2$ for $i<j$ on either side is
$(S_i-S_j)^2/(\alpha S_iS_j)$.  The norm in
\eqref{mat:P-square} includes expectation.
Set $F=\sum_rX_rw_r$ and let $\operatorname{skew}M=(M-M^t)/2$.
Completing the square in the skew-symmetric matrix
$S^{-1/2}[T,b]S^{-1/2}$ bounds the first two terms on the left of
\eqref{mat:commutator-bound} by
\[
 2\alpha\bigl\|\operatorname{skew}(S^{1/2}FS^{1/2})\bigr\|^2.
\]

For completeness, we estimate this square together with the last term.
Write $S_i=su_i$, where $0<u_i\le1$, so
$\alpha S_iB_i=1-u_i$.  The rows of the block matrix
$X=(X_1,\ldots,X_m)$ are orthogonal, since $XX^t=B$ is diagonal.
For column $j$ of the vertically stacked matrix $(w_r)_r$, let $z_{ij}$
be its component along the normalized $i$th row of $X$.  Zero rows can
be completed to orthonormal directions, and cause no contribution to
$F$.  Then $F_{ij}=\sqrt{B_i}\,z_{ij}$ and the remaining column
components are orthogonal to these directions.

For $i<j$, the quadratic form in $(z_{ij},z_{ji})$ has matrix
\[
 \begin{pmatrix}
 1-su_iu_j&-s\sqrt{u_iu_j(1-u_i)(1-u_j)}\\
 -s\sqrt{u_iu_j(1-u_i)(1-u_j)}&1-su_iu_j
 \end{pmatrix}.
\]
Its largest eigenvalue is at most
$1+s(z-2z^2)\le1+s/8$, where $z=\sqrt{u_iu_j}$; here
$(1-u_i)(1-u_j)\le(1-z)^2$.
Diagonal and orthogonal column components have coefficient at most
$1$.  Summation and expectation prove \eqref{mat:commutator-bound}.
\end{proof}

We also use a block version of Cauchy--Schwarz.  If
$X=(X_1,\ldots,X_m)$ and $S^{-1}=\Id/s+\alpha XX^t$, then
\begin{equation}\label{mat:block-bound}
 \left\|a+\sum_rX_rL_r\right\|_S^2
 \le s\|a\|^2+\frac1\alpha\|L\|^2.
\end{equation}
This follows column by column because the block operator
$S^{1/2}(\Id/\sqrt{s},\sqrt{\alpha}X)$ has product with its transpose
equal to $\Id$.

The scalar estimates in Lemma~\ref{sc:bounds} give $p(a)>0$ for every
real $a$.  To reserve part of the negative $H$ term for the cost of
$\log\det S$, fix
\begin{equation}\label{mat:completion-constants}
 \epsilon=\frac3{20},\qquad\kappa=\frac{49}{100},
\end{equation}
and put
\[
 \rho=\epsilon\Id+\kappa S/s.
\]
Then $\Id-\rho>0$.  A weighted completion of the square gives
\begin{align}
 \tau\left[T\left(Yh_0+\sum_rX_rM_r\right)\right]
 &\le\tau\bigl(H(\Id-\rho)p\bigr)
       +\frac{\|M p^{-1/2}\|^2}{4(1-\epsilon)}
       +\frac{\alpha s\|Yh_0p^{-1/2}\|^2}
                    {4(1-\epsilon-\kappa)}.
 \label{mat:weighted-completion}
\end{align}
To verify that the order of the weights is valid, write
$Q_0=H(\Id-\rho)$ and $F=Yh_0+\sum_rX_rM_r$.
Apply $\langle a,b\rangle\le\|a\|^2+\frac14\|b\|^2$ to
$Q_0^{1/2}p^{1/2}$ and $Q_0^{-1/2}TFp^{-1/2}$.
The resulting second square has left weight bounded above by
\[
 T^2Q_0^{-1}\le\alpha S(\Id-\rho)^{-1},
\]
with equality on the range of $H$.  Here inverses are interpreted as
pseudoinverses on the common nullspace of $Q_0$ and $T$.
Since
\[
 (\Id-\rho)S^{-1}
   =\frac{1-\epsilon-\kappa}{s}\Id
           +(1-\epsilon)\alpha XX^t,
\]
the same block argument as in \eqref{mat:block-bound} proves
\eqref{mat:weighted-completion}.  Thus this step does not require $p$
to commute with $S$.

\subsection{Two scalar estimates and the final chaos bound}

The following estimate combines the cost $\tr(T-\log S)$ with the
negative $H$ term left by \eqref{mat:weighted-completion} and the
terms quadratic in $Y$.  A second estimate will control the derivative
residuals, including their dependence on $Z$, by $\mathcal C$ and the
higher-chaos energy.  The proof of the first estimate is given in
Section~\ref{sc:absorption-proof}.  Its strict remainder controls the
eigenvalues of $B$ through the weight \eqref{mat:rigidity-weight}.

\begin{proposition}[Absorption with a first-chaos remainder]
\label{mat:scalar-absorption}
Let $A\in L^2(\gamma_m;\operatorname{Sym}_m)$ satisfy $\E A=0$, and
define $X_r,Y,B$ by \eqref{mat:chaos} and $S,T,H$ by
\eqref{mat:dual-matrices}.  With the constants
\eqref{mat:dual-constants}, \eqref{mat:constants}, and
\eqref{mat:completion-constants}, the functions \eqref{mat:scalar-functions}, and
$\rho=\epsilon\Id+\kappa S/s$, one has
\begin{equation}\label{mat:absorption-inequality}
 \begin{split}
 &\tr(T-\log S)-\tau(H\rho p)
 +s\tau\left[Y^2\left((b-\beta)^2+
       \frac{\alpha h_0^2}{4(1-\epsilon-\kappa)p}\right)\right]\\
 &\hspace{15mm}\le\frac{69}{500}\|Y\|^2
       -\frac1{1000}\sum_{i=1}^m\omega(\lambda_i),
 \end{split}
\end{equation}
where $\lambda_1,\ldots,\lambda_m$ are the eigenvalues of $B$.
All scalar factors in the expected traces are evaluated at $A$.
\end{proposition}

The remaining derivative terms are compared entrywise in an eigenbasis
of $A$.  This comparison retains a multiple of the eigenvalue-separation
surplus $\mathcal C$ rather than discarding it.

\begin{proposition}[Two-eigenvalue estimate]\label{mat:scalar-pair}
Let $a_i,a_j\in\R$ and $G,g,Z\in\R$.  Set
\[
 \delta=\bigl|\log(q(a_i)/q(a_j))\bigr|,\qquad
 r_{ij}=\frac{1+m_0(\min(a_i,a_j))}{16}
       \left(1+\frac{\delta^2}{30}+\frac{\delta^4}{1680}\right),
\]
where $m_0$ is the six-interval lower bound in \eqref{tr:bins}.  Define the endpoint residuals
\[
 \begin{array}{lll}
 D_j=Gb_j+Z\Delta k_{ij},&L_j=Gl_j+Z\Delta b_{ij},&
 M_j=Ge_j+Z\psi_{ij}+cgb_j,\\
 D_i=Gb_i+Z\Delta k_{ji},&L_i=Gl_i+Z\Delta b_{ji},&
 M_i=Ge_i+Z\psi_{ji}+cgb_i.
 \end{array}
\]
Let
$Q(u,v)=\frac{43}{100}u^2+\frac{26}{100}uv+\frac{68}{25}v^2$.
Then
\begin{align}
 &\operatorname{avg}_{v=i,j}\left[
 s(D_v-\beta g)^2+\left(1+\frac s8\right)(D_v-\gamma g)^2
 +\gamma^2(s-1)g^2+\frac{L_v^2}{\alpha}
 +\frac{M_v^2}{4(1-\epsilon)p_v}\right]\notag\\
 &\hspace{15mm}\le Q(G,g-G/2)+r_{ij}\delta^2Z^2.
 \label{mat:pair-inequality}
\end{align}
Here $\operatorname{avg}_{v=i,j}$ is the arithmetic mean of the two
endpoint expressions, counting both also when $i=j$.  Moreover
$p_v=p(a_v)>0$, and all divided differences extend continuously to
$a_i=a_j$.
\end{proposition}

The proof of Proposition~\ref{mat:scalar-pair} occupies
Section~\ref{sc:pair-proof}.  We now show that these two scalar estimates
close the matrix argument.

\begin{proof}[Proof of Theorem~\ref{mat:entropy-smooth}]
Apply Lemma~\ref{mat:expansion} in
\eqref{mat:deficit-reduction}.  Bound its last line by
Lemma~\ref{mat:commutator}, its first square by
\eqref{mat:block-bound}, and its linear $T$ term by
\eqref{mat:weighted-completion}.  The $H$ terms combine to
$-\tau(H\rho p)$.  Since $Y=Y^t$ and $h_0,p,b$ are commuting functions
of $A$, the two terms quadratic in $Y$ have exactly the form appearing
in \eqref{mat:absorption-inequality}, by cyclicity of the trace;
no commutation with $Y$ is required.  Set
$\Omega_B=\sum_{i=1}^m\omega(\lambda_i)$.  Consequently,
\begin{align}
 \mathscr D\le {}&-V_0-\mathcal C+\frac{69}{500}\|Y\|^2
    -\frac{\Omega_B}{1000}
    +\|D-\beta g\|_S^2
    +\left(1+\frac s8\right)\|D-\gamma g\|^2
    +\frac1\alpha\|L\|^2\nonumber\\
 &\hspace{22mm}+\gamma^2\tau\left(T\sum_rg_r^2\right)
    +\frac{\|Mp^{-1/2}\|^2}{4(1-\epsilon)}.
 \label{mat:pre-pair}
\end{align}
Use $S\le s\Id$ and $T\le(s-1)\Id$.  In an eigenbasis of $A$, the
entries of $G_r,g_r,Z_r$ are symmetric in $i,j$.  Pairing the entries
$(i,j)$ and $(j,i)$ therefore gives precisely the endpoint average in
\eqref{mat:pair-inequality}; right multiplication by $p^{-1/2}$ gives
the denominator $p_j$ for entry $(i,j)$.  Proposition~\ref{mat:scalar-pair}
now bounds the last five terms of \eqref{mat:pre-pair} by
\[
 \E\sum_{i,j,r}Q\bigl((G_r)_{ij},(g_r)_{ij}-(G_r)_{ij}/2\bigr)
       +\mathcal C.
\]
The separation surplus is the same $\mathcal C$ as in
Proposition~\ref{tr:surplus}: its coefficient at $(i,j)$ is independent
of the derivative index $r$, so summing over $r$ is unchanged by rotating
that index.  All changes of eigenbasis are pointwise; no eigenvectors are
differentiated.  The separation surplus cancels.
Because $Q$ has constant coefficients,
its summed quadratic form is invariant under orthonormal changes of
basis.  We may therefore return to fixed coordinates and apply the
Gaussian chaos decomposition.

For a degree-$d$ component of $Y$, where $d\ge2$, the associated
components of $g$ and $G$ satisfy $g=G/d$, and
$\|G\|^2=d\|Y\|^2$.  Different degrees remain orthogonal after
differentiation.  The required bound follows, degree by degree, from
\begin{equation}\label{mat:final-polynomial}
 Q(1,u-1/2)+\frac7{50}u\le1-u,
 \qquad 0\le u\le\frac12.
\end{equation}
Indeed, the right side minus the left side factors as
$\frac1{50}(1-2u)(1+68u)\ge0$.  Set $u=1/d$, multiply by
$\|G\|^2$, and sum over the chaoses; the sums converge because
$Y\in H^1(\gamma_m)$.  The summed $Q$ term is at most
$V_0-\frac7{50}\|Y\|^2$.  Substitution into
\eqref{mat:pre-pair} therefore gives
\[
 \mathscr D\le-\frac1{500}\|Y\|^2-\frac{\Omega_B}{1000},
\]
which is \eqref{mat:strict-remainder}.
\end{proof}
